\section{Camera \& Rendering Dynamics}

This section covers the camera spring-damper system and gauge needle animation physics, which use the same mathematical framework as the physics engine.

\cite{ogata2009modern} provides the control theory foundation for spring-damper systems.

\subsection{Camera Spring-Damper (Frequency-Domain Parameterization)}

\begin{equation}
  K = \omega_0^2, \quad C = 2 \zeta \omega_0
  \label{eq:spring-damper-params}
\end{equation}

\textbf{Physical Meaning}: The camera follows the car using a spring-damper system, parameterized by natural frequency $\omega_0$ and damping ratio $\zeta$. This is more intuitive than raw spring/damping constants.

\textbf{Parameters}:
\begin{itemize}
  \item $\omega_0 = 2.2$ rad/s: Natural frequency (default)
  \item $\zeta = 0.8$: Damping ratio (critically damped at $\zeta = 1$)
  \item $\zeta < 1$: Underdamped (slight overshoot)
  \item $\zeta = 1$: Critically damped (no overshoot)
  \item $\zeta > 1$: Overdamped (slow convergence)
\end{itemize}

\subsection{Verlet Camera Integration}

\begin{equation}
  c_{\text{new}} = c + (c - c_{\text{prev}}) \cdot e^{-C \cdot dt} + K \cdot (p_{\text{target}} - c) \cdot dt^2
  \label{eq:verlet-camera}
\end{equation}

\textbf{Physical Meaning}: The camera position is updated using Verlet integration (same method as the car body). The spring force pulls toward the target, while exponential damping represents velocity decay.

\textbf{Source Code}: \coderef{physics.js}{1310--1356}

\begin{lstlisting}[caption={Verlet camera with spring-damper}]
const dampingFactor = Math.exp(-C * dt);
const springForce = K * (targetX - cameraX);
cameraX = cameraX + (cameraX - cameraXPrev) * dampingFactor + springForce * dt * dt;
\end{lstlisting}

\subsection{Camera Zoom}

\begin{equation}
  z_{\text{target}} = z_{\text{max}} - (v - v_{\text{threshold}}) \cdot s \cdot 0.01
  \label{eq:camera-zoom}
\end{equation}

\begin{equation}
  z = z + (z_{\text{target}} - z) \cdot 2 \cdot dt
  \label{eq:zoom-smoothing}
\end{equation}

\textbf{Physical Meaning}: Camera zoom is a linear function of speed: faster speeds zoom out to show more road. The zoom is smoothed via first-order time lag.

\textbf{Parameters}:
\begin{itemize}
  \item $v_{\text{threshold}} = 10$ m/s: Speed at which zoom begins
  \item $s = 0.005$: Zoom sensitivity
  \item $z_{\text{min}}, z_{\text{max}}$: Zoom bounds (typically 0.5 to 2.0)
\end{itemize}

\subsection{Gauge Needle Animation (Spring-Damper)}

\begin{equation}
  F_{\text{spring}} = (x_{\text{target}} - x) \cdot k_s \cdot dt_{\text{norm}}
  \label{eq:needle-spring}
\end{equation}

\begin{equation}
  v = v \cdot c^{dt_{\text{norm}}}
  \label{eq:needle-damping}
\end{equation}

\begin{equation}
  x = x + v
  \label{eq:needle-update}
\end{equation}

\textbf{Physical Meaning}: Gauge needles (tachometer, speedometer, lateral-G meter) animate toward their target values using an asymmetric spring-damper system. Rise and fall have different stiffness values for arcade appeal.

\textbf{Source Code}: \coderef{ui.js}{42--93}

\textbf{Parameters}:
\begin{itemize}
  \item $k_s = 0.070$: Spring stiffness (default)
  \item $c_{\text{rise}} = 0.92$: Damping coefficient for upward motion
  \item $c_{\text{fall}} = 0.88$: Damping coefficient for downward motion
  \item $dt_{\text{norm}} = dt \times 60$: Framerate-independent timestep (normalized to 60 Hz)
\end{itemize}

\textbf{Notes}: Framerate independence is achieved via $dt_{\text{norm}}$, allowing smooth animation at any display refresh rate.

\subsection{Needle Flutter (Vibration Effect)}

\begin{equation}
  x_{\text{flutter}} = x + A(x) \cdot \sin(2\pi f_{\text{flutter}} \cdot t)
  \label{eq:needle-flutter}
\end{equation}

\textbf{Physical Meaning}: High-speed gauge vibration (~1.43 Hz) simulates engine vibration transmission. Amplitude scales with needle position to be more noticeable at high speeds.

\textbf{Parameters}:
\begin{itemize}
  \item $f_{\text{flutter}} \approx 1.43$ Hz: Oscillation frequency
  \item $A = (x - x_{\text{threshold}}) \times 0.012$: Amplitude (position-dependent)
\end{itemize}

\newpage
